% !PRINTED 172
\weberpage{172}
%<<CONT
substitution, these forms are homogeneous functions of
$n$ variables, and hence no more than $n$ mutually
independent ones can exist. This does not mean that all
these forms can be expressed \gesperrt{rationally} in terms of $n$ among them,
but between $n+1$ invariants there must always subsist a rational
equation, which results by elimination of the $n$ variables.

But there remains the converse question to investigate, whether
for a linear substitution group of dimension
$n$ there really always exist $n$ independent invariants.

If we put fixed values for the variables, then all the
invariants of the group thereby likewise receive
determinate values. The question that we still have to answer
is now this: whether to a system of values of the invariants
there exist determinate systems of values of the variables, and
indeed in finite number. Once this is proved, we can regard the
variables as (many-valued) algebraic functions of the invariants,
and the number of mutually independent
invariants cannot be smaller than the number
of the variables, because otherwise a part of the variables, when the
values of the invariants are given, would still remain
arbitrary.

If $(x)$ is a system of values of the variables corresponding
to a determinate system of values of the invariants, and $A$ a substitution
of the group $S$, then by the nature of the invariants
the system $A(x)$ corresponds to the same system of values of the invariants,
and our problem therefore has at least as many
solutions as the degree of the group amounts to. That this, however, is the
exact number of solutions appears from the following
considerations.

In special cases, i.e. for special values of the invariants,
several of these systems of values of the variables
may coincide. But this can happen only for such systems of values
of the $(x)$ for which a relation of the form
$(x) = A(x)$ subsists, for from $A(x) = B(x)$ there would follow $(x) = A^{-1}B(x)$,
which is contained in the form $(x) = A(x)$.

We assume a linear homogeneous function $\Theta$ of the variables
$x_1, x_2, \dots x_n$, which we denote thus:
\weq{1}{\Theta = \Theta(x) = c_1 x_1 + c_2 x_2 + \cdots + c_n x_n.}
